\section{Related Work}
\label{sec:related_work}

This work builds upon three main pillars of research: parameter-efficient expert ensembling, dynamic gating frameworks in non-stationary serving, and the mathematical principles of active inference.

\paragraph{Parameter-Efficient Expert Ensembling:}
Modern multi-expert serving architectures increasingly leverage Parameter-Efficient Fine-Tuning (PEFT) adapters, such as LoRA \cite{lora2021}, to implement specialized task experts without incurring the prohibitive memory footprint of multiple monolithic models \cite{moe2017}. Dynamic ensembling frameworks blend these specialized adapters at test-time using input-dependent ensembling coefficients \cite{lorahub2023, s-lora2024}. However, existing methods typically treat each serving query as independent, ignoring temporal context and rendering them highly susceptible to observation noise and sudden input fluctuations.

\paragraph{Dynamic Gating and Non-Stationary Serving:}
Dynamic model serving in non-stationary environments must balance responsiveness to task switches against stability in stable regimes. Stateless routers like SABLE \cite{sable2024} evaluate queries independently, yielding high responsiveness but suffering from extreme high-frequency routing jitter under noisy inputs. Conversely, stateful gating methods apply low-pass filtering, such as Momentum-Merge \cite{momentummerge2025}, PAC-Kinetics \cite{packinetics2025}, or biochemical ODE ensembling (ChemMerge) \cite{chemmerge2025}. While these stateful methods successfully suppress routing noise, they introduce a fixed temporal lag that severely degrades serving accuracy specifically during rapid, step-by-step task transitions, creating the Jitter-Lag Trade-Off.

\paragraph{Mathematical Principles of Active Inference:}
Active Inference and the Free Energy Principle \cite{friston2006, friston2010} formulate perception and action as a joint variational optimization problem where an agent minimizes its Variational Free Energy to align its internal belief state with its environment. While active inference has achieved outstanding success in theoretical neuroscience and robotic control \cite{millidge2020neural, safron2020active}, its application to deep systems serving has been limited. Existing control-theoretic routing approaches often rely on heuristic step-size scheduling or computationally expensive iterative gradient unrolling. We bridge this gap by deriving an exact, closed-form analytical solver for variational free energy minimization, resolving the systems-level Jitter-Lag Dilemma with microsecond-level serving overhead.
